% P03.1 isolated scalar-defect nonremovability, live integration.

\begin{theorem}
\label{thm:p031-isolated-scalar-defect-nonremovability}
Retain the setting and notation of
\Cref{thm:p031-curvewise-scalar-defect-localization}.  Thus
\[
 \begin{aligned}
 \mathcal B_g&=\{c\in C:h_g(c)<q_g(c)\},\\
 \mathcal F_g&=
 \{c\in C:q_g\text{ is not subharmonic on any neighborhood of }c\}.
 \end{aligned}
\]
Every isolated point of $\mathcal B_g$ belongs to $\mathcal F_g$.
More precisely, if $c$ is isolated in $\mathcal B_g$, then $q_g$
has a strict local maximum at $c$.

Consequently, if $\mathcal B_g$ is closed and discrete, then
\begin{equation}
\label{eq:p031-closed-discrete-scalar-defect-exact-failure}
                         \mathcal F_g=\mathcal B_g,
\end{equation}
and $q_g$ is subharmonic on $C$ if and only if
$\mathcal B_g=\varnothing$.  The same conclusion applies whenever
$\mathcal B_g$ is locally finite.
\end{theorem}

\begin{corollary}
\label{cor:p031-psh-scalar-defect-no-isolated-points}
In the setting of
\Cref{thm:p031-curvewise-scalar-defect-localization}, the existence of an
isolated point of $\mathcal B_g$ implies that $u=s^*$ is not
classically plurisubharmonic on $V$.  Equivalently, if $u$ is
classically plurisubharmonic, then for every analytic extension $K'/K$
and every extended smooth strict test curve $g:C\to V_{K'}$, the strict
canonical-retraction defect $\mathcal B_g$ has no isolated point.  In
particular, a closed discrete or locally finite defect locus must be empty.
\end{corollary}

The theorem is a nonremovability statement, not an existence theorem for
isolated scalar defects.  It neither proves that $\mathcal B_g$ is locally
finite nor excludes a nonempty defect locus without isolated points.
